\documentclass[aps,prl,reprint,10pt,amsmath,amssymb,superscriptaddress,a4paper]{revtex4-2}
\usepackage{graphicx} % Graphics package
\usepackage{dcolumn} % Double column package
\usepackage{amsmath,amsfonts,amsthm} % Math packages
\usepackage[margin=2cm]{geometry} % Sets 2cm margins
\usepackage{datetime} % Package for automatic date & time
\usepackage{subfig}
\usepackage[hidelinks]{hyperref}
\usepackage{pythonhighlight}

\setcitestyle{authoryear,round}

\begin{document}

\title{Holograhy}

\author{T. Nguyen (z5416116)}
\affiliation{UNSW PHYS3117 - Physics Laboratory}
\date{\currenttime~\today}

\begin{abstract}
The aim of this experiment is to verify the past observations of previous experiments and differentiate between the acousto-optic modulator and deflector. We investigated the difference between acousto-optic modulators and deflectors through analysing the Klein-Cook parameter, 
finding that the acoustic width of the crystal in the instrument determined the regime in which the acousto-optic 
effect was operating in. Exploring the various relationships between light and the angle enabled us to develop a more well rounded understanding of the Bragg and Raman-Nath regimes.
\end{abstract}

\maketitle

\section{Introduction}

Photography captures a scene from one specific angle. Through a pinhole, the scene is projected into a black box and the image is 
created when the material in the box is exposed to light, causing a series of chemical reactions \citep{canon}. The image retains 
the information from the scene through diffracting the oncoming light using the pinhole producing a diffracted field pattern. This 
pattern is then focused through a lens onto a local image plane \citep{reece1}. This method of retaining information on a scene 
relies on Fourier optics and the intensity of light reflected from the object. Therefore, the information about the incident light's 
phase is lost.

Holography is a technique for imaging scenes that captures and retains all information of the scene and projects it onto film. The 
principle of holography was discovered by accident by Dennis Gabor in 1947, when he was developing better imaging techniques in 
electron microscopy. On a fateful Easter day, he came up with the idea to take a picture that contains the whole information, and 
then correct it by optical means \citep{gabor}. Realising that in ordinary photography, phase is lost, he added a coherent background, 
or a `reference wave' that he could use to compare with the object wave, to produce fringes on an interference pattern.

Due to the limitations in his era, Gabor ran into a strange effect in his initial experiments \citep{gaboring}. Some parts of the 
object became much brighter than the coherent background at an angle in-line with the imaging plane. The effect looked like a focusing 
of the reference beam onto an image of the object on the opposite side of the film. This is known as the conjugate image and is a 
result of the first order diffraction of the reference beam. This effect can be removed by separating the illuminating beam used on 
the object, and the reference beam, which was done after the invention of the laser in 1961, \citep{leithupatnieks}.

At this point in time, holography had not yet been able to capture the spectacular 3D images we see today. To see the hologram, Gabor 
had to focus the field through a microscope or a short-focus eyepiece, enough to reconstruct the image but the area was not large enough 
to see the image's depth and angular behaviour. It was not until the improvement of the laser to helium-neon as well as bleached or 
`phase' holograms that one could have the impression of peering through a window to look at the object itself.

The first large image capturing took place in 1970 of a statue of Abraham Lincoln \citep{stroke}. Since then, holography had become a 
niche artistic medium. Whether or not it is a relevant creative medium is outside the scope of this report but it is important to note 
its uses, albeit small and perhaps insignificant, in shaping the cultural landscape. There was an attempt that establishing holography 
as a mainstream medium throughout the 1970s and 1980s when artists such as Schweitzer and Moree made holograms in their basement and then 
taught others how to do so. 

The most challenging component is not the interpretation of the art but its implementation and creation involving high powered lasers 
and equipment not available to the average person. With stronger collaborations between artists and scientists may this medium of 
creativity flourish \citep{pepper}.

The culmination of physics, art, medicine and philosophy comes with a remark made by Gabor in his Nobel lecture. He brought up the idea 
that diffuse holograms were a disturbed memory, drawing between the similarities with the noise in creating holograms to the fragmentation 
of human memories. It is an exciting question that gave birth to the Holonomic Brain Theory \citep{pribram}. The model describes the brain 
as a holographic storage network where pieces of long-term memory is distributed over a dendritic arbor. The model has served being a 
foundation for many other quantum models of the brain ever since.

Finally, the last item of note for holography would be the development of The Holographic Principle \citep{hooft}. The core idea is analogous 
to the creation of a holographic photograph where a 3D object is projected onto a 2D plane where all information is retained. In this case, the 
principle concerns itself with a description of particle states in the near vicinity of a black hole.

\section{Theory}

Light is an electromagnetic wave and therefore, its behaviour governed by Maxwell's equations. We could treat our physics this way and continuously 
solve the equations for a set of boundary conditions however it would be tedious and difficult to obtain. Often, optical problems are treated 
with Huygens-Fresnel theory. The link back into Maxwell's electromagnetic theory can be made when considering Kirchhoff's diffraction theory \citep{reece1}.

\subsection{Huygens-Fresnel construction}

Competing with Newton's corpuscle theory of light in the 17th century, was a primitive version of the wave theory of light we see today \citep{huygens}.
Huygens' construction of light involved a series of irregular shock waves with great but finite velocity through the ether. The model we will be using 
is the second part, where each point of the wavefront is itself the origin a secondary spherical wave. Later on, in 1818, following Thomas Young's interference 
experiment, Fresnel showed that Huygens' construction was itself, an interference of waves and therefore could explain the rectilinear propagation of 
light and also diffraction effects \citep{fresnel}.

The optical wave functions can be written simply as 

\begin{equation}
    U(r,t) = Ae^{i(kr-\omega t)}.
\end{equation}

We can see that, despite the lack of true vector representation and Maxwell rigor, our scalar wave function still comprises of an amplitude and a phase. 
The corresponding intensity of the wave function would be the magnitude of this function, 

\begin{equation}
    I(r,t) = U(r,t)U^*(r,t) = |A|^2.
\end{equation}

\subsubsection{Assumptions}



\bibliographystyle{apsrev4-2}
\bibliography{../Notes/references}

\end{document}